\documentclass[10pt,a4paper]{amsart}
\usepackage[margin=19mm]{geometry}
\usepackage{amsmath,amssymb,mathtools}
\usepackage[scr=rsfs,cal=euler]{mathalfa}
\usepackage{xfrac}
\usepackage[unicode,hidelinks,bookmarks=false]{hyperref}
\newcommand{\ie}{i.e.\ }
\newcommand{\cf}{cf.\ }
\newcommand{\resp}{respectively\ }
\newcommand{\Subsection}{\subsection}
\providecommand{\nobreakdash}{\nobreak\hbox{-}}
\setlength{\emergencystretch}{2em}
\multlinegap0pt
\usepackage{bm}
\usepackage{bbm}
\usepackage{dsfont}
\usepackage{xspace}
\usepackage{cancel}
\usepackage{mathtools}
\usepackage{amsthm}
\makeatletter
\@ifundefined{thm}{\newtheorem{thm}{Thm}}{}
\@ifundefined{lem}{\newtheorem{lem}[thm]{Lemma}}{}
\makeatother
\makeatletter
\newcommand{\CC} {\mathcal{C}}
\newcommand{\OO}{\mathcal{O}}
\newcommand{\OOhat}{\widehat{\OO}}
\DeclareMathOperator{\Quot}{Quot}
\newcommand{\bC} {\mathbf{C}}
\newcommand{\bD} {\mathbf{D}}
\newcommand{\beqn}{\begin{equation}}
\newcommand{\eeqn}{\end{equation}}
\expandafter\ifx\csname enumrom\endcsname\relax
\newenvironment{enumrom}
{\renewcommand{\theenumi}{\roman{enumi}} \begin{enumerate}}
{\end{enumerate} \renewcommand{\theenumi}{\arabic{enumi}}}
\fi
\makeatother
\title[Structure theorems for braided Hopf algebras]{Structure theorems for braided Hopf algebras}
\author{Craig Westerland}
\date{Source: arXiv:2406.13874v2 (CC BY 4.0)}
\begin{document}
\maketitle

\noindent\emph{Source and licence.} Structure theorems for braided Hopf algebras. Version arXiv:2406.13874v2 (CC BY 4.0), distributed under the Creative Commons Attribution 4.0 International licence, as recorded in the arXiv licence metadata for this exact version and shown on the abstract page https://arxiv.org/abs/2406.13874v2. Full rights notice: licenses/arxiv-2406.13874-CC-BY-4.0.txt.\par\smallskip

\noindent\textit{Editorial scope.} Counted unit: the Lem whose source label is \texttt{tech\_const\_lem} in the article (source0.tex lines 1454--1458, source label \texttt{tech\_const\_lem}), delivered together with the article's own complete original proof of it (source0.tex lines 1460--1495, one \texttt{proof} environment of 36 lines). No mathematics is written, repaired or re-derived: statement and proof are the article's own text line for line. Required context retained uncounted: none. Every definition, display and label of those lines is retained unchanged. Omitted from this package and not redistributed: 1--1453, 1459--1459, 1496--3262. Nothing inside a retained excerpt is omitted or altered: every retained body is delivered exactly as the article's lines are written, so no omission receipt of any kind accompanies this package. Citations: the retained excerpts carry no citation command at all, so no bibliography entry of the article is transcribed and no reference is redistributed. Reference adaptation: none was needed and none was made; every reference inside the delivered unit resolves inside it, so no unresolved reference or citation is produced. Printed slips kept literal: no printed word, symbol, hypothesis or number of the counted statement or of its proof was altered. Layout and class: the article's own class is \texttt{article} with packages including amsmath, amscd, amsfonts, amsthm, xypic, amssymb, graphicx, mathdots, hyperref, color, cleveref, tikz, pgfplots, braids; the wrapper is the lane's pinned \texttt{amsart} editorial wrapper at 10pt on A4, which restates the theorem-like environments the retained text needs and adds the article's own macro definitions that the retained text needs (\texttt{\textbackslash CC}, \texttt{\textbackslash OO}, \texttt{\textbackslash OOhat}, \texttt{\textbackslash Quot}, \texttt{\textbackslash bC}, \texttt{\textbackslash bD}, \texttt{\textbackslash beqn}, \texttt{\textbackslash eeqn}), with extra wrapper packages (bm, bbm, dsfont, xspace, cancel, mathtools). The delivered numbering is the unit's own. Assets: no retained span contains a figure environment, a graphics-inclusion command, a PDF-page inclusion, a table or a TikZ picture; no image file is copied into this package, the package declares no asset dependency and no image file is redistributed with it. Licence: version arXiv:2406.13874v2 of this article is distributed under the Creative Commons Attribution 4.0 International licence, as recorded in the arXiv licence metadata for this exact version and shown on the abstract page https://arxiv.org/abs/2406.13874v2. A full rights notice is delivered at licenses/arxiv-2406.13874-CC-BY-4.0.txt. Characters: every retained excerpt is pure ASCII, so the delivered engine, XeLaTeX with its default Latin Modern text font, reports no missing character.
\begin{lem} \label{tech_const_lem}

Let $\OO$ be an operad in $\bD$, and $\CC \leq \OOhat$ be an essentially strict braided subsequence of its profinite completion.  Then if $V \in \bC$ is finitely braided, the value of the Artin functor $\CC[V]$ is pro-constant.

\end{lem}

\begin{proof}

Recall that $\OOhat(n)$ is the inverse system of coinvariants
%
$$\{\OO(n)_{\ker_Q}, \; Q \in \Quot_n\},$$
%
where $\ker_Q =\ker(B_n \twoheadrightarrow Q)$.  Since $\CC(n)$ is a subsequence of $\OOhat(n)$, $\CC(n)$ is indexed over the same inverse system.  We will write 
%
$$\CC(n)_{\ker_Q} \leq \OOhat(n)_{\ker_Q}$$
%
for the corresponding term in the inverse system for $\CC(n)$, even though it may not be the $\ker_Q$-coinvariants of some subobject of $\OO(n)$.

If $P \twoheadrightarrow Q$ is a map in $\Quot_n$, write $K^P_Q$ for its kernel.  Then the iterated coinvariants satisfy $(\OOhat(n)_{\ker_P})_{K^P_Q} = \OOhat(n)_{\ker_Q}$.  This yields an inclusion
%
\beqn \label{ess_strict_eqn} (\CC(n)_{\ker_P})_{K^P_Q} \subseteq \CC(n)_{\ker_Q}.\eeqn
%
Since $\CC(n)$ is essentially strict, this inclusion is an equality for sufficiently large $Q$.

Now $\CC_n[V]$ is the pro-object $\{\CC(n)_{\ker_Q} \otimes_{B_n} V^{\otimes n} \}$, again indexed over $Q \in \Quot_n$.  Once $Q$ is large enough that $\ker_Q$ acts trivially on $V^{\otimes n}$, the corresponding term is
%
$$\CC(n)_{\ker_Q} \otimes_{B_n} V^{\otimes n} = \CC(n)_{\ker_Q} \otimes_{Q} V^{\otimes n}.$$
%
If $Q$ is also large enough that (\ref{ess_strict_eqn}) is an equality, we see that for each $P \twoheadrightarrow Q$,
%
$$\CC(n)_{\ker_P} \otimes_{P} V^{\otimes n} = (\CC(n)_{\ker_P})_{K^P_Q} \otimes_Q V^{\otimes n} =  \CC(n)_{\ker_Q} \otimes_{Q} V^{\otimes n},$$
%
and so this system is pro-constant: $\CC_n[V] \cong \CC(n)_{\ker_Q} \otimes_Q V^{\otimes n}$.  

This in turn implies that 
%
$$\CC[V] = \bigoplus_{n=0}^{\infty} \CC_n[V]$$
%
is pro-constant.  Recall that direct sums of pro-objects are indexed on the products of the indexing categories of their terms, so $\CC[V]$ is indexed on $\prod_n \Quot_n$.  For each $n$, we have established that there is some $Q_n \in \Quot_n$ such that the part of the system $\CC_n[V]$ lying over $Q_n$ is constant.  Thus $\CC[V]$ is pro-constant, being constant above $\prod_n Q_n$.


\end{proof}

\end{document}
